\documentclass[11pt]{article}
\usepackage[a4paper,margin=26mm]{geometry}
\usepackage{amsmath,amssymb,mathtools}
\usepackage{unicode-math}
\setmainfont{Latin Modern Roman}
\setmathfont{Latin Modern Math}
\usepackage{xurl,hyperref}
\hypersetup{hidelinks,pdftitle={Problem statement}}
\newcommand{\sref}[2]{\hyperlink{source:#1:#2}{[S#2]}}
\newcommand{\cataloguescope}[1]{#1}
\begin{document}
\section{Reeb components — Question 5.1}\label{statement}

\subsection{Definitions and mathematical statement}

\paragraph{Shared notation.}$\mathbb N=\{1,2,\ldots\}$, and $\mathbb Z,\mathbb Q,\mathbb R,\mathbb C$ denote the integers, rationals, reals and complex numbers. Any other conventions are specified locally.


An open $3$-manifold is noncompact and has no boundary. A Reeb component is a foliated solid torus whose boundary is a torus leaf and whose interior leaves are planes spiraling toward that boundary.
\paragraph{Statement recorded in the supplied source.}
For an open $3$-manifold $M$, determine the lower bounds on the number of Reeb components forced in every codimension-one foliation of $M$, distinguishing zero, finitely many and infinitely many. The source also asks for geometric versions: under uniform bounds on ambient and foliation geometry, constrain Reeb components by bounded diameter or weight them by core-circle length, and in an infinite case study the growth of their necessary number inside metric balls. These bounds and weights are parameters of the geometric formulation, not additional fixed data in the topological question. \sref{statement}{2}
\subsection{Short English statement}
Determine how many Reeb components are unavoidable on open manifolds, and formulate geometric lower bounds and growth versions under uniform geometry constraints.
\subsection{Sources}
\begin{description}
\item[\hypertarget{source:statement:1}{S1}] ULAM.AI and the UnsolvedMath Contributors, \emph{UnsolvedMath: A Curated Collection of Open Mathematics Problems}, catalogue entry. CC BY 4.0. \url{https://huggingface.co/datasets/ulamai/UnsolvedMath}.
\item[\hypertarget{source:statement:2}{S2}] Danny Calegari, \emph{Problems in foliations and laminations of 3-manifolds}, version 0.78 (2002), Question 5.1; Definitions 1.1--1.10 and local definitions and remarks. \url{https://arxiv.org/abs/math/0209081v1}.
\end{description}
\label{end:statement}
\end{document}
